\section{IBP reduction and exact master coefficients}
\label{sec:ibp-coefficients}

\subsection{Cut-decorated topology equivalence}

Many ordered diagram pairs generate families related by loop-momentum
changes.  Reusing an IBP reduction requires more than equality of the
quadratic denominator polynomials.  We call two cut families equivalent when
an affine transformation
\begin{equation}
 \ell_r=\sum_{s=1}^{L}A_{rs}\ell'_s+\sum_aB_{ra}p_a,
 \qquad A_{rs},B_{ra}\in\mathbb Q,
 \qquad |\det A|=1,
 \label{eq:affine-loop-map-main}
\end{equation}
maps their ordered denominators into each other up to a slot permutation,
preserves every propagator power, maps each cut momentum together with its
energy orientation, and preserves the partition into phase-space,
forward-virtual, and conjugate-virtual integration momenta.  It must also
preserve any measured-line label.  Under these conditions the change of
variables proves an exact identity between the powered cut integrals.
Appendix~\ref{app:topology-equivalence} gives the full criterion and an
explicit two-loop example.

This relation is stronger than equality of uncut graphs and weaker than
equality of measured cross-section terms.  The latter also requires the same
numerator projection and measurement assignment.  The distinction prevents a
loop routing or a reversal of a cut momentum from silently changing the
physical integral.

\subsection{Cut-aware IBP reduction}

For a complete family $\mathcal F=\{D_1,\ldots,D_N\}$, IBP identities are
\begin{equation}
 0=\int\prod_{r=1}^{L}\dd^D\ell_r\,
 \frac{\partial}{\partial\ell_i^\mu}
 \left[v^\mu\prod_{j=1}^{N}D_j^{-\nu_j}\right].
 \label{eq:ibp-identity}
\end{equation}
The designated cut slots obey Eq.~\eqref{eq:cut-family-condition} throughout
the reduction; they are not imposed on an uncut result afterward.  Solving
the resulting linear system gives
\begin{equation}
 I_{\mathcal F}(\boldsymbol\nu)
 =\sum_{a=1}^{N_M}
 R_{\boldsymbol\nu a}(s_{ij},D)\,M_a,
 \label{eq:ibp-reduction}
\end{equation}
where every $R_{\boldsymbol\nu a}$ is rational in the invariants and in
$D$.  We use \Kira\ for this exact reduction
\cite{Maierhofer:2017gsa,Klappert:2020nbg}.  Equivalent families are reduced
once, and Eq.~\eqref{eq:affine-loop-map-main} transports the index vector back
to all members of the class.

\subsection{Finite-field reconstruction of master coefficients}

After the IBP rules have been inserted, the coefficient of a master is a sum
of many rational contributions,
\begin{equation}
 C_a(\boldsymbol x,D)
 =\sum_{r}c_{ar}(\boldsymbol x,D),
 \label{eq:master-coefficient-sum}
\end{equation}
where $\boldsymbol x$ denotes a dimensionless invariant basis.  Direct
symbolic addition can generate intermediate expressions far larger than the
final rational function.  Instead, the rational operation graph is recorded
with \Ratracer\ and evaluated over finite fields.  \FireFly\ reconstructs the
multivariate rational function from those modular images
\cite{Magerya:2022bcf,Klappert:2019emp}.  The reconstruction is accepted only
after substitution into an independent prime field and an exact
characteristic-zero comparison wherever the latter is tractable.

No special functions are reconstructed at this stage.  Gamma functions,
branch-sensitive powers, and spin tensors are separated from the rational
operation graph and carried as exact prefactors.  This makes the reconstructed
object precisely the rational coefficient $C_a$, rather than a numerical
fit to the full hard function.

\subsection{Extraction of the partonic coefficient}

The reduced expression is assembled master by master,
\begin{equation}
 \widehat{\mathcal C}^{\mathsf P,(n)}
 =\mathscr F^{\mathsf P}
 \frac{1}{x_ax_bz_h^2}
 \sum_a C_a^{\mathsf P,(n)}(s,t,u,\eps)M_a(s,t,u,\eps).
 \label{eq:post-ibp-hard-coefficient}
\end{equation}
The exact hadronic-to-partonic substitutions in
Eq.~\eqref{eq:hadronic-partonic-map} are applied before the scalar hard
coefficient is declared.  After the distribution product and displayed
fraction monomial are divided out, the result must satisfy
Eq.~\eqref{eq:fraction-independence}.  Residual dependence on a momentum
fraction, an auxiliary light-cone direction, an integration momentum, or an
unreduced scalar product terminates the calculation.  The same criterion is
applied separately to each physical TT azimuthal coefficient.
